% Module 1 section file. Included by ../module1_stellarator_equilibrium_geometry.tex.
\section{Chapter summary}
\label{sec:summary}

\begin{enumerate}[leftmargin=1.6em]
\item Magnetic confinement exploits rapid gyromotion and relatively free motion along magnetic field lines. Closing field lines toroidally avoids simple end loss.
\item A stellarator uses three-dimensional external magnetic geometry to produce field-line twist without relying on a large transformer-driven plasma current.
\item Static ideal-MHD equilibrium satisfies
\begin{equation*}
\nabla p=\mathbf J\times\mathbf B,
\qquad
\nabla\times\mathbf B=\mu_0\mathbf J,
\qquad
\nabla\cdot\mathbf B=0.
\end{equation*}
\item Force balance implies $\mathbf B\cdot\nabla p=0$, so pressure is a flux function when nested surfaces exist.
\item A magnetic surface satisfies $\mathbf B\cdot\nabla s=0$. The normalized toroidal-flux coordinate $s$ labels nested surfaces, while $\rho=\sqrt{s}$ is a convenient radius-like coordinate.
\item Rotational transform $\iota=d\theta/d\zeta$ measures average field-line pitch. The safety factor is $q=1/\iota$ under the project convention.
\item Radial variation of $\iota$ produces magnetic shear. Rational values $\iota=n/m$ identify surfaces important for islands and wave resonances.
\item Straight-field-line coordinates make field lines straight in angular coordinate space. Boozer coordinates are especially useful for particle-orbit and quasisymmetry analysis.
\item Three-dimensional surfaces and field strength are represented efficiently by Fourier series in poloidal and toroidal angles.
\item The vacuum coil field and plasma-current field both contribute to the operating equilibrium.
\item Fixed-boundary calculations prescribe the plasma boundary; free-boundary calculations solve the plasma-vacuum interface with coil fields.
\item A numerical equilibrium is a nonlinear constrained problem requiring radial and spectral convergence studies, force-balance checks, and coordinate-quality diagnostics.
\item Downstream modules require not only $\mathbf B$ but also $B$, derivatives, curvature, metrics, Jacobian, $\iota$, $V'$, profiles, conventions, provenance, and uncertainty information.
\item A physics-informed surrogate must preserve equilibrium constraints and derived-quantity accuracy; it cannot be validated by field-value error alone.
\end{enumerate}

\appendix

\section{Vector identities used in this chapter}
\label{app:vector_identities}

\subsection{Triple-product identities}
For vectors $\mathbf A$, $\mathbf C$, and $\mathbf D$,
\begin{align}
\mathbf A\cdot(\mathbf C\times\mathbf D)
&=\mathbf C\cdot(\mathbf D\times\mathbf A)
=\mathbf D\cdot(\mathbf A\times\mathbf C),\\
\mathbf A\times(\mathbf C\times\mathbf D)
&=\mathbf C(\mathbf A\cdot\mathbf D)
-\mathbf D(\mathbf A\cdot\mathbf C).
\end{align}

\subsection{Magnetic-force identity}
For a differentiable vector field $\mathbf B$,
\begin{equation}
(\nabla\times\mathbf B)\times\mathbf B
=(\mathbf B\cdot\nabla)\mathbf B
-\frac{1}{2}\nabla B^2
-\mathbf B(\nabla\cdot\mathbf B).
\end{equation}
When $\nabla\cdot\mathbf B=0$, this yields Equation~\eqref{eq:magnetic_force_decomposition}.

\subsection{Divergence of a curl}
For a sufficiently smooth vector potential $\mathbf A$,
\begin{equation}
\nabla\cdot(\nabla\times\mathbf A)=0.
\end{equation}
Thus representing $\mathbf B=\nabla\times\mathbf A$ enforces the solenoidal condition analytically.

\section{Derivation of the straight-field-line representation}
\label{app:straight_field_derivation}

Let $\psi$ label magnetic surfaces and let $\theta$ and $\zeta$ be periodic angles. A field tangent to surfaces can be represented using cross products containing $\nabla\psi$:
\begin{equation}
\mathbf B
=A(\psi,\theta,\zeta)\nabla\psi\times\nabla\theta
+C(\psi,\theta,\zeta)\nabla\zeta\times\nabla\psi.
\end{equation}
The two terms are tangent to constant-$\psi$ surfaces. A suitable angular-coordinate choice and flux normalization set $A=1$ and make $C=\iota(\psi)$, giving Equation~\eqref{eq:B_contravariant}. This coordinate construction is nontrivial globally; the equation should be understood as the defining representation of an appropriate straight-field-line system, not as a statement that arbitrary geometric angles already have this property.

\section{Example project interface checklist}
\label{app:interface_checklist}

Before a Module 1 state product is accepted by another module, verify:
\begin{enumerate}[leftmargin=1.6em]
\item The equilibrium identifier and source checksum are present.
\item SI units or normalization constants are present for every dimensional field.
\item The meanings of $s$, $\rho$, $\theta$, and $\zeta$ are documented.
\item Angle orientation, flux signs, and the $q$--$\iota$ convention are documented.
\item The valid domain and last closed flux surface are identified.
\item The field-period convention and Fourier convention are identified.
\item $\mathbf B$, $B$, position, Jacobian, metric, and $\iota$ are mutually consistent.
\item Derivatives are generated from the same representation as the primary field where possible.
\item The minimum Jacobian is positive under the stored orientation.
\item Force, divergence, tangency, and transform residuals pass specified tolerances.
\item Radial and spectral convergence evidence is attached.
\item Interpolation and extrapolation behavior is documented.
\item The geometry does not silently extrapolate outside the valid plasma or vacuum domain.
\item Downstream outputs retain the equilibrium identifier.
\end{enumerate}

\section{Review questions and exercises}
\label{app:exercises}

\begin{enumerate}[leftmargin=1.8em]
\item Explain why a purely toroidal magnetic field is not generally sufficient for good toroidal confinement. Distinguish particle streaming from cross-field drifts.
\item Starting from $\nabla p=\mathbf J\times\mathbf B$, prove that both $\mathbf B$ and $\mathbf J$ are tangent to constant-pressure surfaces.
\item Derive Equation~\eqref{eq:Jperp} for the perpendicular current.
\item Use the magnetic-force identity to interpret magnetic pressure and tension.
\item For $B=3\ \mathrm T$ and $p=1.5\times10^5\ \mathrm{Pa}$, calculate the local beta.
\item A field line has $\iota=0.42$. How many poloidal turns does it make during ten toroidal turns? Is the line guaranteed to close?
\item For $\iota(s)=0.3+0.4s$, locate the surfaces with $\iota=1/2$ and $\iota=2/3$, if they lie in $0\leq s\leq1$.
\item Show that Equation~\eqref{eq:B_contravariant} gives $d\theta/d\zeta=\iota$.
\item Explain why $\rho=\sqrt{s}$ should not automatically be interpreted as geometric distance from the magnetic axis.
\item Describe the difference between fixed-boundary and free-boundary equilibrium.
\item Explain why a vacuum coil field is not necessarily the same as a finite-beta plasma equilibrium.
\item What physics is excluded when globally nested surfaces are imposed?
\item Why can accurate interpolation of $B$ still produce inaccurate guiding-center orbits if $\nabla B$ is noisy?
\item Write the Fourier phase for a mode with poloidal number $m=3$, toroidal number $n=2$, and $N_{\mathrm{fp}}=4$ under the convention in Equation~\eqref{eq:R_fourier}.
\item Explain the distinction between stellarator symmetry and quasisymmetry.
\item Give three quantities from Module 1 that Module 4 needs and explain why.
\item Construct a convergence table for $\iota(0.5)$ as $m_{\max}$, $n_{\max}$, and $N_s$ are increased. What pattern would support numerical convergence?
\item Propose a divergence-free neural representation for $\mathbf B$ and identify one disadvantage.
\item Explain why a low mean-square error in $B$ does not guarantee a useful equilibrium surrogate.
\item State the most important assumptions that must accompany any claim based on this Module 1 model.
\end{enumerate}

